\answer{ex:06:1}{$e=(1-e^{-T_a/\tau_e})e^{-d/\tau_e}$. (a)~$0.110$ contra $4.5\cdot10^{-5}$, $\approx2.5\cdot10^3$ vezes. (b)~$\tau_e^*=T_a/\ln(1+T_a/d)=0.59\,\text{s}$.}
\answer{ex:06:2}{Taxa $\nu_x\nu_y(A_+\tau_+-A_-\tau_-)=0.8A_+$ por segundo, $\approx2900A_+$ por hora; $\tau_-=\tau_+/0.6=33\ms$.}
\answer{ex:06:3}{$\mathbf e_1$ se $\mu^2+s_1^2>s_2^2$; aqui $1.01>0.25$, e o neurônio aprende $\mathbf e_1$, embora a dispersão ao longo de $\mathbf e_2$ seja $25$ vezes maior: a regra encontra o vetor principal da matriz de correlação, não da de covariância.}
\answer{ex:06:4}{$V=Re^{-(t_c+L-t)/\tau_\gamma}$ entre a luz e o suco, de modo que $\dot V=V/\tau_\gamma$; pico de área $Re^{-L/\tau_\gamma}$: $0.61R$ e $0.78R$.}
\answer{ex:06:5}{Relação sinal-ruído $1/(N+1)$, tentativas $\propto N+1$: $5\cdot10^5$ tentativas $\approx1$~mês e $5\cdot10^{11}$ tentativas $\approx8\cdot10^4$~anos.}
\answer{ex:06:6}{(a)~$k_2=1/\tau_d$, $k_1=1/\tau_d-1/\tau_\gamma$. (b)~Erro espúrio $-(V/\tau_\gamma)\,\varepsilon/(1+\varepsilon)$, $\varepsilon=\tau_d/\tau_\gamma$, donde $\tau_d\le0.105\,\text{s}$.}
\answer{ex:06:7}{$0.46$, $0.90$, $0.99$, $0.95$ e $0.70$ com $d=3\,\text{s}$. Os pontos caem numa única curva da fração do traço $F=(1-e^{-T_{\text{cue}}/\tau_e})e^{-d/\tau_e}$: aprendizado com $F\gtrsim0.1$, escolha aleatória com $F<10^{-2}$.}
\answer{ex:06:8}{$\eta^*=1/\bigl(2\sigma^2(N+2)\bigr)$, em $N+2$ tentativas; com $m$ flutuando de $N$, em $N(m+2)/m\ge N+2$: a esparsidade não ajuda.}
